\chapter{Body Dysmorphic Disorder}
\index{Body Dysmorphic Disorder}

\section*{Definition and Overview}
Body dysmorphic disorder (BDD) is a mental health condition in which a person is preoccupied with a perceived flaw or defect in their appearance that is either minor or not visible to others. The person spends excessive time and mental energy on this preoccupation, which causes significant distress and interferes with daily functioning. BDD is not vanity -- it is a genuine and distressing illness that is strongly linked to low mood and, in severe cases, to thoughts of self-harm, and it responds well to appropriate treatment. Although the perceived flaw is usually in the person's own mind, the suffering it causes is very real.

\section*{Epidemiology}
Body dysmorphic disorder affects roughly 1--2\% of the population, making it more common than is often realised. It affects men and women fairly equally, although the features they focus on tend to differ. BDD usually begins in adolescence, a time when appearance and self-image are particularly sensitive. Many people with BDD are embarrassed by their symptoms and do not tell anyone, so the condition is under-recognised. It commonly occurs alongside depression, anxiety, and social anxiety disorder.

\section*{Aetiology and Pathophysiology}
The causes of BDD are not fully understood, but several factors appear to contribute:

\begin{itemize}
    \item \textbf{Brain processing differences:} people with BDD tend to focus on details of their appearance rather than the whole, and to process faces and images differently
    \item \textbf{Genetic factors:} BDD is more common in people with a family history of the condition or of obsessive-compulsive disorder
    \item \textbf{Perfectionism and low self-esteem:} a tendency to be self-critical and to place great value on appearance
    \item \textbf{Life experiences:} bullying, teasing, or criticism about appearance, often during childhood or adolescence
    \item \textbf{Social and cultural pressure:} media and social media emphasis on idealised and often edited images of appearance
    \item \textbf{The vicious cycle:} repeated checking, comparing, and avoidance keeps the preoccupation and distress going
\end{itemize}

\section*{Clinical Features}
The key feature is an intense preoccupation with appearance that is out of proportion to reality:

\begin{itemize}
    \item Constant worry about a specific feature, commonly the skin, hair, nose, or weight
    \item Repetitive behaviours: checking mirrors, seeking reassurance, comparing with others, picking at the skin, or camouflaging with make-up or clothing
    \item Avoidance of social situations, photographs, or being seen in certain lighting
    \item Difficulty concentrating at work, school, or in relationships
    \item Distress, shame, and low mood about the perceived flaw
    \item In some people, a mistaken belief that the flaw is real and obvious to everyone (a delusional belief)
    \item Some people seek repeated cosmetic treatments for a feature others cannot see as flawed
\end{itemize}

\section*{Differential Diagnosis}
Other conditions can look similar to BDD:

\begin{itemize}
    \item Normal appearance concerns, which cause less distress and do not dominate life
    \item Obsessive-compulsive disorder, although BDD focuses specifically on appearance
    \item Eating disorders such as anorexia nervosa, where weight and shape concerns are central
    \item Depression and social anxiety disorder, which often coexist with BDD
    \item Gender-related concerns about appearance
    \item Skin picking disorder or other repetitive behaviour disorders
\end{itemize}

\section*{When to Seek Medical Care}
See a doctor or mental health professional if:

\begin{itemize}
    \item Appearance concerns take up a large part of your day or cause significant distress
    \item Checking, camouflaging, or avoidance is affecting work, study, or relationships
    \item You are thinking about repeated cosmetic procedures for a feature others cannot see as flawed
    \item The concerns are accompanied by low mood, anxiety, or social withdrawal
    \item Symptoms have lasted more than a few weeks or are getting worse
\end{itemize}

\textbf{Seek immediate help if there are thoughts of suicide or self-harm.} Contact local emergency services, a local crisis service, or the nearest hospital emergency department. Acute mental-health crises are medical emergencies.

\section*{Diagnosis and Investigations}
\begin{itemize}
    \item \textbf{Clinical interview:} a careful discussion of the appearance concerns, how long they have been present, and how they affect daily life
    \item \textbf{Screening questions:} asking about checking behaviour, avoidance, and the impact of appearance concerns
    \item \textbf{Assessment of mood:} screening for depression, anxiety, and obsessive-compulsive symptoms, which commonly coexist
    \item \textbf{Risk assessment:} asking about thoughts of self-harm or suicide is an essential part of the assessment
    \item \textbf{Ruling out other causes:} in rare cases, physical conditions can affect appearance, and these are considered during assessment
    \item \textbf{No routine tests:} BDD is diagnosed on the basis of the clinical interview; blood tests and scans are not usually needed
\end{itemize}

\section*{Management}
BDD is treatable, and the most effective treatments are psychological and medical:

\begin{itemize}
    \item \textbf{Cognitive behaviour therapy (CBT):} helps people recognise and challenge unhelpful thoughts about appearance and reduce checking, comparing, and avoidance
    \item \textbf{Antidepressant medication:} selective serotonin reuptake inhibitors (SSRIs), prescribed and monitored by a doctor, are effective, often in combination with CBT
    \item \textbf{Treating co-existing conditions:} managing the depression, anxiety, or obsessive-compulsive symptoms that often accompany BDD
    \item \textbf{Involving family:} family members can learn how to support without reinforcing the illness
    \item \textbf{Advice on cosmetic treatment:} cosmetic procedures rarely help BDD and may make the preoccupation worse; this should be discussed openly
    \item \textbf{Specialist referral:} to a psychiatrist or psychologist experienced in BDD for more severe or complex cases
    \item \textbf{Crisis care:} hospital or community crisis services for people at risk of self-harm
\end{itemize}

\section*{Complications}
\begin{itemize}
    \item Depression, social anxiety, and difficulty functioning at work, school, or socially
    \item Withdrawal from relationships and activities
    \item Repeated or unnecessary cosmetic procedures, sometimes with poor results
    \item Skin damage from repetitive skin picking
    \item Thoughts of suicide and self-harm -- the most serious complication
    \item Substance use as a way of coping with distress
    \item Long-term suffering and reduced quality of life if the condition is not treated
\end{itemize}

\section*{Prognosis and Outcomes}
With treatment, many people with BDD improve substantially, although the condition often runs a long course and may flare up in times of stress. CBT and SSRIs help most people reduce the distress and the time spent on the preoccupation. Early treatment is associated with better outcomes, while untreated BDD tends to persist. Some people need longer-term or repeated courses of treatment, and the risk of self-harm means ongoing monitoring is important, especially in severe cases.

\section*{Prevention and Self-Care}
\begin{itemize}
    \item Seek help early, when the concerns first begin to interfere with life
    \item Limit social media and comparison with heavily edited images of appearance
    \item Practise redirecting attention away from checking and comparing
    \item Build self-esteem through activities and relationships that are not about appearance
    \item Tell a trusted person about the struggle, rather than keeping it secret
    \item Follow the treatment plan and attend follow-up appointments
    \item Learn to recognise early warning signs of worsening symptoms and seek help promptly
\end{itemize}

\section*{Sources and Further Reading}
The following sources provide reliable, up-to-date information on body dysmorphic disorder:

\begin{itemize}
    \item NHS. \textit{NHS guidance on body dysmorphic disorder.} https://www.nhs.uk/conditions/
    \item Mayo Clinic. \textit{Information on body dysmorphic disorder and treatment.} https://www.mayoclinic.org/
    \item MSD Manual. \textit{Overview of body dysmorphic disorder.} https://www.msdmanuals.com/
    \item MedlinePlus. \textit{Body dysmorphic disorder patient information.} https://medlineplus.gov/
    \item World Health Organization. \textit{Mental health and disorder information.} https://www.who.int/
    \item NICE. \textit{Obsessive-compulsive and body dysmorphic disorder guidelines.} https://www.nice.org.uk/
\end{itemize}
